\documentclass[11pt]{article}

\usepackage[T1]{fontenc}
\usepackage[utf8]{inputenc}
\usepackage{lmodern}
\usepackage[margin=1in]{geometry}
\usepackage{amsmath}
\usepackage{booktabs}
\usepackage{hyperref}
\usepackage{url}
\usepackage{microtype}
\sloppy

\hypersetup{
  colorlinks=true,
  linkcolor=black,
  urlcolor=blue,
  citecolor=black
}

\title{Deterministic Audit of LLM-Generated Research Claims}
\author{Vladimir V.~Borovikov\\
\small Independent researcher\\
\small \url{https://orcid.org/0009-0005-7773-3581}}
\date{October 2026}

\begin{document}
\maketitle

\begin{abstract}
Large language models are increasingly used to draft the narrative parts of
research: reports, verdicts, discussion sections. The failure mode that matters
is not obvious fabrication but \emph{plausible overstatement}: the number is
right, the interpretation is not, and the prose is confident. We describe a
protocol that makes such overstatement mechanically detectable, and report what
happened when we applied it to our own corpus of 1{,}416 probes across ten
experimental series.

The protocol rests on six requirements: pre-registration freeze, external memory
instead of context recall, separation of generator and auditor, deterministic
parity checking of every reported number, a three-tier taxonomy of claim errors,
and multi-model adversarial review. We contribute an executable checklist of
twelve items and report its honest status on our own work: six pass, one is
amended, one fails, two are partial, two require human judgement. The protocol
is not fully applied, and the checklist says so.
\end{abstract}

\section{Introduction}

When an LLM writes up an experiment, the output is fluent by construction. The
fluency is not evidence of correctness, but it reads as such---to reviewers, and
to the authors. In our own series we observed a repeated pattern: the first
draft of a verdict claimed more than the archive supported, and the claim was
withdrawn only after a separate recount.

Examples from our logs, all verifiable in the public repository:

\begin{itemize}
  \item a null result ($p = 0.338$) described as a refutation ``with 100\%
        strictness'';
  \item a comparison of $2/72$ against $0/72$ reported as a ``law''
        (Fisher $p = 0.4965$~\cite{fisher1935});
  \item a $p$-value printed as $4.81\times10^{-13}$ where the exact two-sided
        Fisher test gives $2.909\times10^{-18}$---a tool error, not a model
        error, that survived six weeks;
  \item a footnote pointing to a data file that contained none of the cited
        bytes.
\end{itemize}

The last two are instructive. The third was produced by our own checker; the
fourth by a model writing a confident sentence about a file it had not opened.
Neither was caught by re-reading.

\section{A taxonomy of claim errors}
\label{sec:taxonomy}

We distinguish three tiers, because they are caught by different means.

\textbf{Tier 0 --- reference fabrication.} The text cites an object that does
not exist: a file, a line, a dataset, a value. Numbers may be arithmetically
perfect. Neither recomputation nor proofreading detects it, because there is
nothing to recompute. Detection requires checking existence.

\textbf{Tier I --- interpretive overstatement.} The archive supports the number,
not the sentence. ``$p = 0.34$'' becomes ``refuted''. Repaired by
reformulation.

\textbf{Tier II --- numerical error.} The archive holds a different number than
the text. Repaired only by recounting from the log.

Tier 0 is the most dangerous and the least defended: tiers I and II have
mechanical checks, tier 0 has only attention.

\section{The protocol}

\paragraph{P1. Freeze.} Code, prompts, hypotheses and endpoints are hashed
before the first probe. A self-signed manifest is a declaration, not a freeze:
without a trusted timestamp it proves only the current state of the files. We
introduce the status \textsc{freeze-amended} for series whose files changed
after the freeze but where the change is documented, the results were recomputed
and matched, and regression tests passed. It differs from ``annulled'' in that
the numbers survive, and from ``frozen'' in that byte-level pipeline
reproduction is lost.

\paragraph{P2. External memory.} Numbers are read from files, never recalled
from context. A model that must remember will eventually remember wrongly.

\paragraph{P3. Separation of contours.} The generator writes, the auditor
counts. The auditor is independent only if it is written from a separate
specification; being ``a different script'' is not independence.

\paragraph{P4. Parity check.} Every reported number---$p$-values, fractions,
percentages, confidence bounds---is recomputed by a script, and the script's
output is published alongside the paper. Two additions proved essential in
practice. First, a \emph{known-answer test}: a deliberately corrupted reference
that the checker must reject, otherwise a blind checker passes everything.
Second, statistics must not be hand-rolled without cross-validation; a parity
check that compares one homemade implementation with another compares a bug with
a bug. Our own Fisher implementation carried an absolute tolerance of
$10^{-12}$, which at $p \approx 10^{-18}$ is $3.4\times10^{5}$ times larger than
the value itself and inflated the tail by five orders of magnitude. It was found
by an independent implementation, not by the protocol---which is why
cross-validation is now a required step rather than a lucky accident.

\paragraph{P5. Taxonomy.} As in Section~\ref{sec:taxonomy}, with Tier~0 checked
by a script that extracts file paths and line ranges from a document and
verifies that the targets exist.

\paragraph{P6. Multi-model review.} Several model families review each other. We
report this as the weakest pillar: the number of passes is a working convention,
not a demonstrated optimum, and models trained on overlapping corpora err
correlatedly, so an ``ensemble'' does not guarantee independence of errors.

\section{Application and honest status}

The checklist has twelve items and is executable. On our corpus it reports:

\begin{center}
\begin{tabular}{@{}lrl@{}}
\toprule
Status & Count & Meaning \\
\midrule
Pass    & 6 & checked mechanically \\
Amended & 1 & post-freeze edit, documented and recomputed \\
Fail    & 1 & verifier independence: same author wrote both \\
Partial & 2 & pre-registration, verdict versioning \\
Manual  & 2 & trusted timestamp, archive immutability \\
\bottomrule
\end{tabular}
\end{center}

The single failure is not repairable by code. It requires a second author
writing the checker from a specification they did not write themselves. We state
this rather than paper over it, because a checklist that can only be passed by
misreporting is worse than no checklist.

The coverage of P4 is ten $p$-values---nine two-sample contrasts and the
post-hoc five-bin test---twenty-one fractions, and all reported
percentages including Clopper--Pearson bounds~\cite{clopper1934}, recomputed by
two independent scripts whose agreement is verified on every run.

\section{Limitations}

The audit sample is not random: we audited the series where an audit was
performed. We make no claim about the series we did not audit.

The nine hypotheses were reconstructed from the manuscript, not frozen before
the run. Multiple-comparison correction (Holm)~\cite{holm1979} is discussed prospectively; it
was not applied to the completed series.

Provenance is not provable: without a trusted timestamp we cannot show that the
archive was in its present state \emph{before} the run rather than adjusted
after. What is demonstrated is narrower---the numbers reproduce from the archive
as it exists in the repository today.

\section{Availability}

Code, probes and the executable checklist:
\url{https://github.com/skywhiteangels-cell/llm-token-architect}.
The accompanying study: \url{https://doi.org/10.5281/zenodo.23679671}.

\paragraph{Model identity.} All probe series reported here ran on
\texttt{flash\_lite}, a publicly available model tier; the verdict texts were
written with a different general-purpose model, noted in the project archive.
Nothing is withheld behind a proprietary claim or a non-disclosure agreement:
the archive and the verification scripts are released under the MIT licence, so
every number can be recomputed without knowing which model produced it.

\begin{thebibliography}{9}

\bibitem{clopper1934}
C.~J. Clopper and E.~S. Pearson.
\newblock The use of confidence or fiducial limits illustrated in the case of
  the binomial.
\newblock \emph{Biometrika}, 26(4):404--413, 1934.

\bibitem{fisher1935}
R.~A. Fisher.
\newblock \emph{The Design of Experiments}.
\newblock Oliver \& Boyd, Edinburgh, 1935.

\bibitem{holm1979}
S.~Holm.
\newblock A simple sequentially rejective multiple test procedure.
\newblock \emph{Scandinavian Journal of Statistics}, 6(2):65--70, 1979.

\end{thebibliography}

\end{document}
